% ============ Chapter 2. Prime ideals (原书 pp.20-45) ============
\chapter{Prime ideals}\label{ch:2}

The notion of prime ideal is central to commutative ring theory. The set
$\Spec A$ of prime ideals of a ring $A$ is a topological space, and the
`localisation' of rings and modules with respect to this topology is an
important technique for studying them. These notions are discussed in \S\ref{sec:4}.
Starting with the topology of $\Spec A$, we can define the dimension of $A$ and
the height of a prime ideal as notions with natural geometrical content. In
\S\ref{sec:5} we treat elementary dimension theory using only field theory, developing
especially the dimension theory of ideals in polynomial rings, including the
Hilbert Nullstellensatz. We also discuss, as an example of an application of the
notion of dimension, the theory of Forster and Swan on estimates for the
number of generators of a module. (Dimension theory will be the subject of
a detailed study in Chapter~\ref{ch:5} using methods of ring theory). In \S\ref{sec:6} we discuss
the classical theory of primary decomposition as modernised by Bourbaki.

% ============ §4 Localisation and Spec of a ring (原书 pp.20-29) ============
\section{Localisation and Spec of a ring}\label{sec:4}\index{Localisation, section 4}\index{Spec of a ring, section 4}

Let $A$ be a ring and $S \subset A$ a multiplicative set; that is (as in
\S\ref{sec:1}) suppose that
\[
  \text{(i) } x, y \in S \Rightarrow xy \in S, \quad \text{and} \quad \text{(ii) } 1 \in S.
\]

\begin{definition}
Suppose that $f\colon A \to B$ is a ring homomorphism satisfying
the two conditions
\begin{enumerate}[(1)]
  \item $f(x)$ is a unit of $B$ for all $x \in S$;
  \item if $g\colon A \to C$ is a homomorphism of rings taking every element of $S$ to
        a unit of $C$ then there exists a unique homomorphism
        $h\colon B \to C$ such that $g = hf$;
\end{enumerate}
then $B$ is uniquely determined up to isomorphism, and is called the
\emph{localisation}\index{localisation} or the \emph{ring of fractions} of $A$ with respect to $S$. We write
$B = S^{-1}A$ or $A_S$, and call $f\colon A \to A_S$ the canonical map.
\end{definition}

We prove the existence of $B$ as follows: define the relation $\sim$ on the
set $A \times S$ by
\[
  (a, s) \sim (b, s') \iff \exists t \in S \text{ such that } t(s'a - sb) = 0;
\]
it is easy to check that this is an equivalence relation (if we just had
$s'a = sb$ in the definition, the transitive law would fail when $S$ has
zero-divisors). Write $a/s$ for the equivalence class of $(a, s)$ under $\sim$, and
let $B$ be the set of these; sums and products are defined in $B$ by the usual
rules for calculating with fractions:
\[
  a/s + b/s' = (as' + bs)/ss', \qquad (a/s) \cdot (b/s') = ab/ss'.
\]
This makes $B$ into a ring, and defining $f\colon A \to B$ by $f(a) = a/1$ we see that
$f$ is a homomorphism of rings satisfying (1) and (2) above. Indeed, if $s \in S$
then $f(s) = s/1$ has the inverse $1/s$; and if $g\colon A \to C$ is as in (2) then we just
have to set $h(a/s) = g(a)g(s)^{-1}$ (the reader should check that $a/s = b/s'$
implies $g(a)g(s)^{-1} = g(b)g(s')^{-1}$). From this construction we see that the
kernel of the canonical map $f\colon A \to A_S$ is given by
\[
  \Ker f = \{a \in A \mid sa = 0 \text{ for some } s \in S\}.
\]
Hence $f$ is injective if and only if $S$ does not contain any zero-divisors
of $A$. In particular, the set of all non-zero-divisors of $A$ is a multiplicative
set; the ring of fractions with respect to $S$ is called the \emph{total ring of fractions}
of $A$. If $A$ is an integral domain then its total ring of fractions is the same
thing as its field of fractions.

In general, let $f\colon A \to B$ be any ring homomorphism, $I$ an ideal of $A$
and $J$ an ideal of $B$. According to the conventions at the beginning of the
book, we write $IB$ for the ideal $f(I)B$ of $B$. This is called the \emph{extension} of
$I$ to $B$, or the \emph{extended ideal}\index{extended ideal}, and is sometimes also written $I^e$. Moreover,
we write $J \cap A$ for the ideal $f^{-1}(J)$ of $A$. This is called the \emph{contracted ideal}\index{contracted ideal}
of $J$, and is sometimes also written $J^c$. In this notation, the inclusions
\[
  I^{ec} \supset I \quad \text{and} \quad J^{ce} \subset J
\]
follow immediately from the definitions; from the first inclusion we get
$I \subset I^e$, but substituting $J = I^e$ in the second gives $I^{ece} \subset I^e$, and hence
\[
  (*) \quad I^{ece} = I^e, \quad \text{and similarly} \quad J^{cec} = J^c.
\]
This shows that there is a canonical bijection between the sets $\{IB \mid I$
is an ideal of $A\}$ and $\{J \cap A \mid J$ is an ideal of $B\}$.

If $P$ is a prime ideal of $B$ then $B/P$ is an integral domain, and since
$A/P^c$ can be viewed as a subring of $B/P$ it is also an integral domain, so
that $P^c$ is a prime ideal of $A$. (The extended ideal of a prime ideal does
not have to be prime.)

An ideal $J$ of $B$ is said to be \emph{primary}\index{primary ideal} if it satisfies the two conditions:
(1) $1 \notin J$, and (2) for $x, y \in B$, if $xy \in J$ and $x \notin J$ then $y^n \in J$ for some
$n > 0$; in other words, all zero-divisors of $B/J$ are nilpotent. The property
that all zero-divisors are nilpotent passes to subrings, so that just as for
prime ideals we see that the contraction of a primary ideal remains primary.
If $J$ is primary then $\sqrt{J}$ is a prime ideal (see Ex.~\ref{ex:4.1-ex}).

The importance of rings of fractions for ring theory stems mainly from
the following theorem.

\begin{theorem}\label{thm:4.1}
\begin{enumerate}[(i)]
  \item All the ideals of $A_S$ are of the form $IA_S$, with $I$ an ideal of $A$.
  \item Every prime ideal of $A_S$ is of the form $pA_S$ with $p$ a prime ideal of
        $A$ disjoint from $S$, and conversely, $pA_S$ is prime in $A_S$ for every such $p$;
        exactly the same holds for primary ideals.
\end{enumerate}
\end{theorem}

\begin{proof}
(i) If $J$ is an ideal of $A_S$, set $I = J \cap A$. If $x = a/s \in J$ then
$x \cdot f(s) = f(a) \in J$, so that $a \in I$, and then $x = (1/s) \cdot f(a) \in IA_S$. The converse
inclusion $IA_S \subset J$ is obvious, so that $J = IA_S$.

(ii) If $P$ is a prime ideal of $A_S$ and we set $p = P \cap A$, then $p$ is a prime
ideal of $A$, and from the above proof $P = pA_S$. Moreover, since $P$ does
not contain units of $A_S$, we have $p \cap S = \emptyset$. Conversely, if $p$ is a prime
ideal of $A$ disjoint from $S$ then
\[
  \frac{a}{s} \cdot \frac{b}{t} \in pA_S \quad \text{with } s, t \in S \Rightarrow rab \in p \text{ for some } r \in S,
\]
and since $r \notin p$ we must have $a \in p$ or $b \in p$, so that $a/s$ or $b/t \in pA_S$. One
also sees easily that $1 \notin pA_S$, so that $pA_S$ is a prime ideal of $A_S$.

For primary ideals the argument is exactly the same: if $p$ is a primary
ideal of $A$ disjoint from $S$ and if $rab \in p$ with $r \in S$, then since no power
of $r$ is in $p$ we have $ab \in p$. From this we get either $a/s \in pA_S$ or $(b/t)^n \in pA_S$
for some $n$.
\end{proof}

\begin{corollary*}
If $A$ is Noetherian (or Artinian) then so is $A_S$.
\end{corollary*}

\begin{proof}
This follows from (i) of the theorem.
\end{proof}

We now give examples of rings of fractions $A_S$ for various multiplicative
sets $S$.

\begin{example}\label{ex:4.1}
Let $a \in A$ be an element which is not nilpotent, and set
$S = \{1, a, a^2, \ldots\}$. In this case we sometimes write $A_a$ for $A_S$. (The reason
for not allowing $a$ to be nilpotent is so that $0 \notin S$. In general if $0 \in S$ then
from the construction of $A_S$ it is clear that $A_S = 0$, which is not very
interesting.) The prime ideals of $A_a$ correspond bijectively with the prime
ideals of $A$ not containing $a$.
\end{example}

\begin{example}\label{ex:4.2}
Let $p$ be a prime ideal of $A$, and set $S = A - p$. In this case
we usually write $A_p$ for $A_S$. (Writing $A_p$ and $A_{A-p}$ to denote the same thing
is totally illogical notation, and the Bourbaki school avoids $A_S$, writing
$S^{-1}A$ instead; however, the notation $A_S$ does not lead to any confusion.)
The localisation $A_p$ is a local ring with maximal ideal $pA_p$. Indeed, as we saw
in Theorem 1, $pA_p$ is a prime ideal of $A_p$, and furthermore, if $J \subset A_p$ is any
proper ideal then $I = J \cap A$ is an ideal of $A$ disjoint from $A - p$, and so $I \subset p$,
giving $J = IA_p \subset pA_p$. The prime ideals of $A_p$ correspond bijectively with
the prime ideals of $A$ contained in $p$.
\end{example}

\begin{example}\label{ex:4.3}
Let $I$ be a proper ideal of $A$ and set $S = 1 + I = \{1 + x \mid
x \in I\}$. Then $S$ is a multiplicative set, and the prime ideals of $A_S$
correspond bijectively with the prime ideals $p$ of $A$ such that $I + p \neq A$.
\end{example}

\begin{example}\label{ex:4.4}
Let $S$ be a multiplicative set, and set $\bar{S} = \{a \in A \mid ab \in S \text{ for
some } b \in A\}$. Then $\bar{S}$ is also a multiplicative set, called the \emph{saturation}\index{saturation (of a multiplicative set)}
of $S$. Since quite generally a divisor of a unit is again a unit, we see from
the definition of the ring of fractions that $A_S = A_{\bar{S}}$, and $\bar{S}$ is maximal among
multiplicative sets $T$ such that $A_S = A_T$. Indeed, one sees easily that $\bar{S} =
\{a \in A \mid a/1 \text{ is a unit in } A_S\}$. The multiplicative set $S = A - p$ of
Example 2 is already saturated.
\end{example}

\begin{theorem}\label{thm:4.2}
Localisation commutes with passing to quotients by ideals.
More precisely, let $A$ be a ring, $S \subset A$ a multiplicative set, $I$ an ideal of $A$
and $\bar{S}$ the image of $S$ in $A/I$; then
\[
  A_S/IA_S \cong (A/I)_{\bar{S}}.
\]
\end{theorem}

\begin{proof}
Both sides have the universal property for ring homomorphisms
$g\colon A \to C$ such that
\begin{enumerate}[(1)]
  \item every element of $S$ maps to a unit of $C$,
  \item every element of $I$ maps to $0$;
\end{enumerate}
the isomorphism follows by the uniqueness of the solution to a universal
mapping problem. In concrete terms the isomorphism is given by
\[
  a/s \bmod IA_S \mapsto \bar{a}/\bar{s}, \quad \text{where } \bar{a} = a + I, \; \bar{s} = s + I.
\]
\end{proof}

In particular, if $p$ is a prime ideal of $A$ then
\[
  A_p/pA_p \cong (A/p)_{A-p}.
\]
The left-hand side is the \emph{residue field} of the local ring $A_p$, whereas the
right-hand side is the field of fractions of the integral domain $A/p$. This
field is written $\kappa(p)$ and called the \emph{residue field of $p$}.\index{residue field of a prime ideal}

\begin{theorem}\label{thm:4.3}
Let $A$ be a ring, $S \subset A$ a multiplicative set, and
$f\colon A \to A_S$ the canonical map. If $B$ is a ring, with ring homomorphisms
$g\colon A \to B$ and $h\colon B \to A_S$ satisfying
\begin{enumerate}[(1)]
  \item $f = hg$,
  \item for every $b \in B$ there exists $s \in S$ such that $g(s) \cdot b \in g(A)$,
\end{enumerate}
then $A_S$ can also be regarded as a ring of fractions of $B$. More precisely,
\[
  A_S = B_{g(S)} = B_T, \quad \text{where } T = \{t \in B \mid h(t) \text{ is a unit of } A_S\}.
\]
\end{theorem}

\begin{proof}
We can factorise $h$ as $B \to B_T \to A_S$; write $\alpha\colon B_T \to A_S$ for
the second of these maps. Now $g(S) \subset T$, so that the composite $A \to B
\to B_T$ factorises as $A \to A_S \to B_T$; write $\beta\colon A_S \to B_T$ for the second of
these maps. Then
\[
  \alpha\beta = 1 \quad \text{and} \quad \beta\alpha = 1,
\]
so that $\beta\alpha = 1$, the identity map of $A_S$. Moreover by assumption, for $b \in B$
there exist $a \in A$ and $s \in S$ such that $bg(s) = g(a)$. Hence, $\beta(a/s) = g(a)/g(s) =
b/1$. In particular for $t \in T$, if we take $a \in A$ such that $t/1 = \beta(a)$ then
$a = \alpha\beta(a) = \alpha(t/1) = h(t)$, so that $a$ is a unit of $A_S$. Hence, $b/t = \beta(a/s)\beta(a^{-1})$,
and $\beta$ is surjective. Thus $\alpha$ and $\beta$ are mutually inverse, giving an
isomorphism $A_S \cong B_T$. The fact that $A_S \cong B_{g(S)}$ can be proved similarly.
(Alternatively, this follows since $T$ is the saturation of the multiplicative
set $g(S)$. The reader should check this for himself.)
\end{proof}

\begin{corollarynum}\label{cor:4.1}
If $p$ is a prime ideal of $A$, $S = A - p$ and $B$ satisfies the
conditions of the theorem, then setting $P = pA_S \cap B$ we have $A_S = B_P$.
\end{corollarynum}

\begin{proof}
Under these circumstances the $T$ in the theorem is exactly $B - P$.
\end{proof}

\begin{corollarynum}\label{cor:4.2}
Let $S \subset A$ be a multiplicative set not containing any zero-
divisors; then $A$ can be viewed as a subring of $A_S$, and for any intermediate
ring $A \subset B \subset A_S$, the ring $A_S$ is a ring of fractions of $B$.
\end{corollarynum}

\begin{corollarynum}\label{cor:4.3}
If $S$ and $T$ are two multiplicative subsets of $A$ with $S \subset T$,
then writing $\bar{T}$ for the image of $T$ in $A_S$, we have $(A_S)_{\bar{T}} = A_T$.
\end{corollarynum}

\begin{corollarynum}\label{cor:4.4}
If $S \subset A$ is a multiplicative set and $P$ is a prime ideal of $A$
disjoint from $S$ then $(A_S)_{PA_S} = A_P$. In particular if $P \subset Q$ are prime ideals of
$A$, then
\[
  (A_P)_{QA_P} = A_Q.
\]
\end{corollarynum}

\begin{definition}
The set of all prime ideals of a ring $A$ is called the \emph{spectrum}\index{spectrum of a ring}
of $A$, and written $\Spec A$; the set of maximal ideals of $A$ is called the
\emph{maximal spectrum} of $A$, and written m-Spec $A$.
\end{definition}

By Theorem~\ref{thm:1.1} we have
\[
  A \neq 0 \iff \text{m-Spec } A \neq \emptyset \iff \Spec A \neq \emptyset.
\]
If $I$ is any ideal of $A$, we set
\[
  V(I) = \{p \in \Spec A \mid p \supset I\}.
\]
Then
\[
  V(I) \cup V(I') = V(I \cap I') = V(II'),
\]
and for any family $\{I_\lambda\}_{\lambda \in \Lambda}$ of ideals we have
\[
  \bigcap_{\lambda \in \Lambda} V(I_\lambda) = V\left(\sum_{\lambda \in \Lambda} I_\lambda\right).
\]
From this it follows that $\mathcal{F} = \{V(I) \mid I \text{ is an ideal of } A\}$ is closed under
finite unions and arbitrary intersections, so that there is a topology on
$\Spec A$ for which $\mathcal{F}$ is the set of closed sets. This is called the \emph{Zariski
topology}\index{Zariski topology}. From now on we will usually consider the spectrum of a ring
together with its Zariski topology. m-Spec $A$ will be considered with the
subspace topology, which we will also call the Zariski topology.

For $a \in A$ we set $D(a) = \{p \in \Spec A \mid a \notin p\}$; this is the complement
of $V(aA)$, and so is an open set. Conversely, any open set of $\Spec A$ can
be written as the union of open sets of the form $D(a)$. Indeed, if
$U = \Spec A - V(I)$ then $U = \bigcup_{a \in I} D(a)$. Hence, the open sets of the form
$D(a)$ form a basis for the topology of $\Spec A$.

Let $f\colon A \to B$ be a ring homomorphism. For $P \in \Spec B$, the ideal
$P \cap A = f^{-1}(P)$ is a point of $\Spec A$. The map $\Spec B \to \Spec A$ defined by
taking $P$ into $P \cap A$ is written ${}^a f$. As one sees easily, $({}^a f)^{-1}(V(I)) = V(IB)$,
so that ${}^a f$ is continuous. If $g\colon B \to C$ is another ring homomorphism then
obviously ${}^a(gf) = {}^a f \cdot {}^a g$. Hence, the correspondence $A \mapsto \Spec A$ and
$f \mapsto {}^a f$ defines a contravariant functor from the category of rings to the
category of topological spaces. If ${}^a f(P) = p$, that is if $P \cap A = p$, we say
that $P$ \emph{lies over} $p$.\index{lies over}

\begin{remark}
For $P$ a maximal ideal of $B$ it does not necessarily follow that
$P \cap A$ is a maximal ideal of $A$; for an example we need only consider the
natural inclusion $A \to B$ of an integral domain $A$ in its field of fractions
$B$. Thus the correspondence $A \mapsto$ m-Spec $A$ is not functorial. This is one
reason for thinking of $\Spec A$ as more important than m-Spec $A$. On the
other hand, one could say that $\Spec A$ contains too many points; for
example, the set $\{p\}$ consisting of a single point is closed in $\Spec A$ if and
only if $p$ is a maximal ideal (in general the closure of $\{p\}$ coincides with
$V(p)$), so that $\Spec A$ almost never satisfies the separation axiom $T_1$.
\end{remark}

Let $M$ be an $A$-module and $S \subset A$ a multiplicative set; we define the
localisation $M_S$ of $M$ in the same way as $A_S$. That is,
\[
  M_S = \left\{\frac{m}{s} \;\middle|\; m \in M, \; s \in S\right\},
\]
and
\[
  \frac{m}{s} = \frac{m'}{s'} \iff t(s'm - sm') = 0 \text{ for some } t \in S.
\]
If we define addition in $M_S$ and scalar multiplication by elements
of $A_S$ by
\[
  m/s + m'/s' = (s'm + sm')/ss' \quad \text{and} \quad (a/s) \cdot (m/s') = am/ss',
\]
then $M_S$ becomes an $A_S$-module, and a canonical $A$-linear map $M \to M_S$
is given by $m \mapsto m/1$; the kernel is $\{m \in M \mid sm = 0 \text{ for some } s \in S\}$. If $S =
A - p$ is the complement of a prime ideal $p$ of $A$ we write $M_p$ for $M_S$. The set
$\{p \in \Spec A \mid M_p \neq 0\}$ is called the \emph{support} of $M$,\index{support of a module} and written $\Supp(M)$. If $M$ is
finitely generated, and we let $M = A\omega_1 + \cdots + A\omega_n$, then
\[
  p \in \Supp(M) \iff M_p \neq 0 \iff \exists i \text{ such that } \omega_i \neq 0 \text{ in } M_p
\]
\[
  \iff \exists i \text{ such that } \ann(\omega_i) \subset p \iff \ann(M) = \bigcap_{i=1}^{n} \ann(\omega_i) \subset p,
\]
so that $\Supp(M)$ coincides with the closed subset $V(\ann(M))$ of $\Spec A$.

\begin{theorem}\label{thm:4.4}
$M_S \cong M \otimes_A A_S$.
\end{theorem}

\begin{proof}
The map $M \times A_S \to M_S$ defined by $(m, a/s) \mapsto am/s$ is $A$-bilinear,
so that there exists a linear map $\alpha\colon M \otimes A_S \to M_S$ such that $\alpha(m \otimes a/s) =
am/s$. Conversely we can define $\beta\colon M_S \to M \otimes A_S$ by $\beta(m/s) = m \otimes (1/s)$;
indeed, if $m/s = m'/s'$ then $ts'm = tsm'$ for some $t \in S$, and so
\[
  m \otimes (1/s) = m \otimes (ts'/tss') = ts'm \otimes (1/tss') = tsm' \otimes (1/tss')
                 = m' \otimes (1/s').
\]
Now it is easy to check that $\alpha$ and $\beta$ are mutually inverse $A_S$-linear
maps. Hence, $M_S$ and $M \otimes_A A_S$ are isomorphic as $A_S$-modules.
\end{proof}

\begin{theorem}\label{thm:4.5}
$M \mapsto M_S$ is an exact (covariant) functor from the category
of $A$-modules to the category of $A_S$-modules. That is, for a morphism of
$A$-modules $f\colon M \to N$ there is a morphism $f_S\colon M_S \to N_S$ of $A_S$-modules
such that
\[
  (1_M)_S = 1_{M_S}, \qquad (gf)_S = g_S f_S,
\]
and such that an exact sequence $0 \to M' \to M \to M'' \to 0$ goes into an
exact sequence $0 \to M'_S \to M_S \to M''_S \to 0$.
\end{theorem}

\begin{proof}
To prove the exactness of $0 \to M'_S \to M_S$ on the last line, view $M'$
as a submodule of $M$; then for $x \in M'$ and $s \in S$,
\[
  x/s = 0 \text{ in } M_S \iff tx = 0 \text{ for some } t \in S
\]
\[
  \iff x/s = 0 \text{ in } M'_S,
\]
as required. The remaining assertions follow from the properties of the
tensor product (see Appendix A) and from the previous theorem. (Of
course they can easily be proved directly.)
\end{proof}

It follows from this that localisation commutes with $\otimes$ and with $\Tor$,
and we will treat all this together in the section on flatness in \S\ref{sec:7}.

Let $A$ be a ring, $M$ an $A$-module and $p \in \Spec A$. There are at least
two interpretations of what it should mean that some property of $A$ or $M$
holds `locally at $p$'. Namely, this could mean that $A_p$ (or $M_p$) has the
property, or it could mean that $A_q$ (or $M_q$) has the property for all $q$ in
some neighbourhood $U$ of $p$ in $\Spec A$. The first of these is more commonly
used, but there are cases when the two interpretations coincide. In any
case, we now prove a number of theorems which assert that a local property
implies a global one.

\begin{theorem}\label{thm:4.6}
Let $A$ be a ring, $M$ an $A$-module and $x \in M$. If $x = 0$ in $M_p$
for every maximal ideal $p$ of $A$, then $x = 0$.
\end{theorem}

\begin{proof}
\[
  x = 0 \text{ in } M_p \iff sx = 0 \text{ for some } s \in A - p \iff \ann(x) \not\subset p.
\]
However, if $1 \notin \ann(x)$ then by Theorem~\ref{thm:1.1}, there must exist a maximal
ideal containing $\ann(x)$. Therefore $1 \in \ann(x)$, that is $x = 0$.
\end{proof}

\begin{theorem}\label{thm:4.7}
Let $A$ be an integral domain with field of fractions $K$; set
$\mathcal{X} = \text{m-Spec } A$. We consider any ring of fractions of $A$ as a subring of $K$.
Then in this sense we have
\[
  A = \bigcap_{m \in \mathcal{X}} A_m.
\]
\end{theorem}

\begin{proof}
For $x \in K$ the set $I = \{a \in A \mid ax \in A\}$ is an ideal of $A$. Now
$x \in A_m$ is equivalent to $I \not\subset m$, so that if $x \in A_m$ for every maximal ideal
$m$ then $1 \in I$, that is $x \in A$.
\end{proof}

\begin{remark}
The above $I$ is the ideal consisting of all possible denominators
of $x$ when written as a fraction of elements of $A$, together with $0$, and this
can be called the \emph{ideal of denominators} of $x$; similarly $Ix$ can be called
the \emph{ideal of numerators} of $x$.
\end{remark}

\begin{theorem}\label{thm:4.8}
Let $A$ be a ring and $M$ a finite $A$-module. If $M \otimes_A \kappa(m) = 0$
for every maximal ideal $m$ then $M = 0$.
\end{theorem}

\begin{proof}
$\kappa(m) = A_m/mA_m$, so that $M \otimes_A \kappa(m) = M_m/mM_m$, and by NAK
(Theorem~\ref{thm:2.2}), $M \otimes_A \kappa(m) = 0 \iff M_m = 0$. Thus the assertion follows from
Theorem~\ref{thm:4.6}.
\end{proof}

The theorem just proved is easy, but we can weaken the assumption
that $M$ is finite over $A$; we have the following result.

\begin{theorem}\label{thm:4.9}
Let $f\colon A \to B$ be a homomorphism of rings, and $M$ a finite
$B$-module; if $M \otimes_A \kappa(p) = 0$ for every $p \in \Spec A$, then $M = 0$.
\end{theorem}

\begin{proof}
If $M \neq 0$ then by Theorem 6 there is a maximal ideal $P$ of $B$ such
that $M_P \neq 0$, so that by NAK, $M_P/PM_P \neq 0$. If we now set $p = P \cap A$ then,
since $pM_P \subset PM_P$, we have $M_P/pM_P \neq 0$. Set $T = B - P$ and $S = A - p$;
then the localisation $M_S = M_p$ of $M$ as an $A$-module and the localisation
$M_T$ of $M$ as a $B$-module coincide (both of them are $\{m/s \mid m \in M, \; s \in S\}$).
We have $f(S) \subset T$, so that
\[
  M_p = M_T \neq 0 \quad \text{and hence} \quad M_p/PM_p = (M_p/PM_p)_T = (M \otimes_A \kappa(p))_T \neq 0;
\]
it follows that $M \otimes_A \kappa(p) \neq 0$.
\end{proof}

\begin{remark}
In Theorem~\ref{thm:4.9} we cannot restrict $p$ to be a maximal ideal of
$A$. As one sees from the proof we have just given, $M = 0$ provided that
$M \otimes \kappa(p) = 0$ for every ideal $p$ which is the restriction of a maximal ideal
of $B$. However, if for example $(A, m)$ is a local integral domain with field
of fractions $B$, and $M = B$, then $M \otimes \kappa(m) = B/mB = 0$, but $M \neq 0$.
\end{remark}

\begin{theorem}\label{thm:4.10}
Let $A$ be a ring and $M$ a finite $A$-module.
\begin{enumerate}[(i)]
  \item For any non-negative integer $r$ set
        \[
          U_r = \{p \in \Spec A \mid M_p \text{ can be generated over } A_p \text{ by } r \text{ elements}\};
        \]
        then $U_r$ is an open subset of $\Spec A$.
  \item If $M$ is a module of finite presentation then the set
        \[
          U = \{p \in \Spec A \mid M_p \text{ is a free } A_p\text{-module}\}
        \]
        is open in $\Spec A$.
\end{enumerate}
\end{theorem}

\begin{proof}
(i) Suppose that $M_p = A_p\omega_1 + \cdots + A_p\omega_r$. Each $\omega_i$ is of the form
$\omega_i = m_i/s_i$ with $s_i \in A - p$ and $m_i \in M$, but since $s_i$ is a unit of $A_p$, we can
replace $\omega_i$ by $m_i$, and so assume that $\omega_i$ is (the image in $M_p$ of) an element
of $M$. Define a linear map $\varphi\colon A^r \to M$ from the direct sum of $r$ copies of
$A$ to $M$ by $(a_1, \ldots, a_r) \mapsto \sum a_i\omega_i$, and write $C$ for the cokernel of $\varphi$. Localising
the exact sequence $A^r \to M \to C \to 0$ at a prime ideal $q$, we get an
exact sequence
\[
  A_q^r \to M_q \to C_q \to 0,
\]
and when $q = p$ we get $C_p = 0$. $C$ is a quotient of $M$, so is finitely generated,
so that the support $\Supp(C)$ is a closed set, and hence there is an open
neighbourhood $V$ of $p$ such that $C_q = 0$ for $q \in V$. This means that $V \subset U_r$.
(In short, if $\omega_1, \ldots, \omega_r \in M$ generate $M_p$ at $p$ then they generate $M_q$ for all $q$
in a neighbourhood of $p$.)

(ii) Suppose that $M_p$ is a free $A_p$-module, and let $\omega_1, \ldots, \omega_r$ be a basis.
As in (i) there is no loss of generality in assuming that $\omega_i \in M$. Moreover,
if we choose a suitable $D(a)$ as a neighbourhood of $p$ in $\Spec A$, $\omega_1, \ldots, \omega_r$
generate $M_q$ for every $q \in D(a)$. Thus, replacing $A$ by $A_a$ and $M$ by $M_a$,
we can assume that the elements $\omega_1, \ldots, \omega_r$ satisfy $M_q = \sum A_q\omega_i$ for every
prime ideal $q$ of $A$. Then by Theorem 6,
\[
  M/\sum A\omega_i = 0, \quad \text{that is } M = A\omega_1 + \cdots + A\omega_r.
\]
(We think of replacing $A$ by $A_a$ as shrinking $\Spec A$ down to the
neighbourhood $D(a)$ of $p$.) Now, defining $\varphi\colon A^r \to M$ as above, and letting
$K$ be its kernel, we get the exact sequence
\[
  0 \to K \to A^r \to M \to 0;
\]
moreover, $K_p = 0$. By Theorem~\ref{thm:2.6}, $K$ is finitely generated, so that applying
(i) with $r = 0$, we have that $K_q = 0$ for every $q$ in a neighbourhood $V$ of
$p$; this gives $(A_q)^r \cong M_q$, so that $V \subset U$.
\end{proof}

% ============ 习题 ============
\section*{Exercises to \S\ref{sec:4}}
\emph{Prove the following propositions.}

\begin{exercise}\label{ex:4.1-ex}
The radical of a primary ideal is prime; also, if $I$ is a proper ideal containing
a power $m^\nu$ of a maximal ideal $m$ then $I$ is primary and $\sqrt{I} = m$.
\end{exercise}

\begin{exercise}\label{ex:4.2-ex}
If $P$ is a prime ideal of a ring $A$ then the \emph{symbolic $n$th power} of $P$ is the ideal
$P^{(n)}$ given by
\[
  P^{(n)} = P^nA_P \cap A.
\]
This is a primary ideal with radical $P$.
\end{exercise}

\begin{exercise}\label{ex:4.3-ex}
If $S$ is a multiplicative set of a ring $A$ then $\Spec(A_S)$ is homeomorphic to
the subspace $\{p \mid p \cap S = \emptyset\} \subset \Spec A$; this is in general neither open nor
closed in $\Spec A$.
\end{exercise}

\begin{exercise}\label{ex:4.4-ex}
If $I$ is an ideal of $A$ then $\Spec(A/I)$ is homeomorphic to the closed subset
$V(I)$ of $\Spec A$.
\end{exercise}

\begin{exercise}\label{ex:4.5-ex}
The spectrum of a ring $\Spec A$ is quasi-compact, that is, given an open
covering $\{U_\lambda\}_{\lambda \in \Lambda}$ of $X = \Spec A$ (with $X = \bigcup_\lambda U_\lambda$), a finite number of the
$U_\lambda$ already cover $X$.
\end{exercise}

\begin{exercise}\label{ex:4.6-ex}
If $\Spec A$ is disconnected then $A$ contains an idempotent $e$ (an element $e$
satisfying $e^2 = e$) distinct from $0$ and $1$.
\end{exercise}

\begin{exercise}\label{ex:4.7-ex}
If $A$ and $B$ are rings then $\Spec(A \times B)$ can be identified with the disjoint
union $\Spec A \sqcup \Spec B$, with both of these open and closed in
$\Spec(A \times B)$.
\end{exercise}

\begin{exercise}\label{ex:4.8-ex}
If $M$ is an $A$-module, $N$ and $N'$ submodules of $M$, and $S \subset A$ a
multiplicative set, then $N_S \cap N'_S = (N \cap N')_S$, where both sides are consi-
dered as subsets of $M_S$.
\end{exercise}

\begin{exercise}\label{ex:4.9-ex}
A topological space is said to be \emph{Noetherian} if the closed sets satisfy the
descending chain condition. If $A$ is a Noetherian ring then $\Spec A$ is a
Noetherian topological space. (Note that the converse is not true in
general.)
\end{exercise}

\begin{exercise}\label{ex:4.10-ex}
We say that a non-empty closed set $V$ in a topological space is \emph{reducible} if
it can be expressed as a union $V = V_1 \cup V_2$ of two strictly smaller closed
sets $V_1$ and $V_2$, and \emph{irreducible} if it does not have any such expression. If
$p \in \Spec A$ then $V(p)$ is an irreducible closed set, and conversely every
irreducible closed set of $\Spec A$ can be written as $V(p)$ for some $p \in \Spec A$.
\end{exercise}

\begin{exercise}\label{ex:4.11-ex}
Any closed subset of a Noetherian topological space can be written as a
union of finitely many irreducible closed sets.
\end{exercise}

\begin{exercise}\label{ex:4.12-ex}
Use the results of the previous two exercises to prove the following: for $I$ a
proper ideal of a Noetherian ring, the set of prime ideals containing $I$ has
only finitely many minimal elements.
\end{exercise}

% ============ §5 The Hilbert Nullstellensatz and first steps in dimension theory (原书 pp.30-37) ============
\section{The Hilbert Nullstellensatz and first steps in dimension theory}\label{sec:5}\index{Hilbert Nullstellensatz, section 5}\index{dimension theory, section 5}

Let $X$ be a topological space; we consider strictly decreasing (or
strictly increasing) chains $Z_0 \supsetneq Z_1 \supsetneq \cdots \supsetneq Z_r$ of length $r$ of irreducible closed
subsets of $X$. The supremum of the lengths, taken over all such chains, is
called the \emph{combinatorial dimension}\index{combinatorial dimension} of $X$ and denoted $\dim X$. If $X$ is a
Noetherian space then there are no infinite strictly decreasing chains, but it
can nevertheless happen that $\dim X = \infty$.

Let $Y$ be a subspace of $X$. If $S \subset Y$ is an irreducible closed subset of $Y$ then
its closure in $X$ is an irreducible closed subset $S' \subset X$ such that $S' \cap Y = S$.
Indeed, if $S' = V \cup W$ with $V$ and $W$ closed in $X$ then $S = (V \cap Y) \cup (W \cap Y)$,
so that say $S = V \cap Y$, but then $V = S'$. It follows easily from this that
$\dim Y \le \dim X$.

Let $A$ be a ring. The supremum of the lengths $r$, taken over all strictly
decreasing chains $p_0 \supsetneq p_1 \supsetneq \cdots \supsetneq p_r$ of prime ideals of $A$, is called the \emph{Krull
dimension}\index{Krull dimension}, or simply the \emph{dimension} of $A$, and denoted $\dim A$. As one sees
easily from Ex.~\ref{ex:4.10-ex}, the Krull dimension of $A$ is the same thing as the
combinatorial dimension of $\Spec A$. For a prime ideal $p$ of $A$, the supremum
of the lengths, taken over all strictly decreasing chains of prime ideals
$p = p_0 \supsetneq p_1 \supsetneq \cdots \supsetneq p_r$ starting from $p$, is called the \emph{height} of $p$,\index{height of a prime ideal} and denoted
$\htop p$; (if $A$ is Noetherian it will be proved in Theorem~\ref{thm:13.5} that $\htop p < \infty$).
Moreover, the supremum of the lengths, taken over all strictly increasing
chains of prime ideals $p = p_0 \subsetneq p_1 \subsetneq \cdots \subsetneq p_r$ starting from $p$, is called the
\emph{coheight} of $p$, and written $\coht p$. It follows from the definitions that
\[
  \htop p = \dim A_p, \quad \coht p = \dim A/p \quad \text{and} \quad \htop p + \coht p \le \dim A.
\]

\begin{remark}
In more old-fashioned terminology $\htop p$ was usually called the
\emph{rank} of $p$, and $\coht p$ the \emph{dimension} of $p$; in addition, Nagata [N1] calls
$\dim A$ the \emph{altitude} of $A$.
\end{remark}

\begin{example}\label{ex:5.1}
The prime ideals in the ring $\Z$ of rational integers are the ideals
$p\Z$ generated by the primes $p = 2, 3, 5, \ldots$, together with $(0)$. Hence, every $p\Z$
is a maximal ideal, and $\dim \Z = 1$. More generally, any principal ideal
domain which is not a field is one-dimensional.
\end{example}

\begin{example}\label{ex:5.2}
An Artinian ring is zero-dimensional; indeed, we have seen in
the proof of Theorem~\ref{thm:3.2} that there are only a finite number of maximal
ideals $p_1, \ldots, p_n$, and that the product of all of these is nilpotent. If then $p$ is a
prime ideal, $p \supset (0) = (p_1 \cdots p_n)^\nu$ so that $p \supset p_i$ for some $i$; hence, $p = p_i$, so
that every prime ideal is maximal.
\end{example}

\begin{example}\label{ex:5.3}
A zero-dimensional integral domain is just a field.
\end{example}

\begin{example}\label{ex:5.4}
The polynomial ring $k[X_1, \ldots, X_n]$ over a field $k$ is an integral
domain, and since
\[
  k[X_1, \ldots, X_n]/(X_1, \ldots, X_i) \cong k[X_{i+1}, \ldots, X_n],
\]
$(X_1, \ldots, X_i)$ is a prime ideal of $k[X_1, \ldots, X_n]$. Thus
\[
  (0) \subset (X_1) \subset (X_1, X_2) \subset \cdots \subset (X_1, \ldots, X_n)
\]
is a chain of prime ideals of length $n$, and $\dim k[X_1, \ldots, X_n] \ge n$. In fact we
will shortly be proving that equality holds.
\end{example}

For an ideal $I$ of a ring $A$ we define the \emph{height of $I$} to be the infimum of the
heights of prime ideals containing $I$:
\[
  \htop I = \inf\{\htop p \mid I \subset p \in \Spec A\}.
\]
Here also we have the inequality
\[
  \htop I + \dim A/I \le \dim A.
\]
If $M$ is an $A$-module we define the \emph{dimension of $M$} by
\[
  \dim M = \dim(A/\ann(M)).
\]
If $M$ is finitely generated then $\dim M$ is the combinatorial dimension of the
closed subspace $\Supp(M) = V(\ann(M))$ of $\Spec A$.

A strictly increasing (or decreasing) chain $p_0, p_1, \ldots$ of prime ideals is said
to be \emph{saturated}\index{saturated chain} if there do not exist prime ideals strictly contained between
any two consecutive terms. We say that $A$ is a \emph{catenary ring}\index{catenary ring} if the following
condition is satisfied: for any prime ideals $p$ and $p'$ of $A$ with $p \subset p'$, there
exists a saturated chain of prime ideals starting from $p$ and ending at $p'$, and
all such chains have the same (finite) length.

If a local domain $(A, m)$ is catenary then for any prime ideal $p$ we have
$\htop p + \coht p = \dim A$. Conversely, if $A$ is a Noetherian local domain and
this equality holds for all $p$ then $A$ is catenary (Ratliff [3], 1972); the proof
of this is difficult, and we postpone it to Theorem 31.4. Practically all the
important Noetherian rings arising in applications are known to be
catenary; the first example of a non-catenary Noetherian ring was
discovered in 1956 by Nagata [15].

We now spend some time discussing the elementary theory of dimensions
of rings which are finitely generated over a field $k$.

\begin{theorem}\label{thm:5.1}
Let $k$ be a field, $L$ an algebraic extension of $k$ and
$a_1, \ldots, a_n \in L$; then
\begin{enumerate}[(i)]
  \item $k[a_1, \ldots, a_n] = k(a_1, \ldots, a_n)$.
  \item Write $\varphi\colon k[X_1, \ldots, X_n] \to k[a_1, \ldots, a_n]$ for the homomorphism
        over $k$ which maps $X_i$ to $a_i$; then $\Ker \varphi$ is the maximal ideal generated
        by $n$ elements of the form $f_1(X_1), f_2(X_1, X_2), \ldots, f_n(X_1, \ldots, X_n)$, where
        each $f_i$ can be taken to be monic in $X_i$.
\end{enumerate}
\end{theorem}

\begin{proof}
Let $f_1(X)$ be the minimal polynomial of $a_1$ over $k$; then $(f_1(X_1))$ is a
maximal ideal of $k[X_1]$, so that $k[a_1] \cong k[X_1]/(f_1(X_1))$ is a field, and hence
$k[a_1] = k(a_1)$. Now let $\varphi_2(X)$ be the minimal polynomial of $a_2$ over $k(a_1)$;
then since $k(a_1) = k[a_1]$, the coefficients of $\varphi_2$ can be expressed as
polynomials in $a_1$, and there is a polynomial $f_2 \in k[X_1, X_2]$, monic in $X_2$,
such that $\varphi_2(X_2) = f_2(a_1, X_2)$. Thus
\[
  k[a_1, a_2] = k[a_1][a_2] = k(a_1)[a_2] \cong k(a_1)[X_2]/(f_2(a_1, X_2)).
\]
Proceeding in the same way, for $1 \le i \le n$ there is an
$f_i(X_1, \ldots, X_i) \in k[X_1, \ldots, X_i]$, monic in $X_i$, such that
\[
  k[a_1, \ldots, a_i] = k(a_1, \ldots, a_i)
\]
\[
  \cong k(a_1, \ldots, a_{i-1})[X_i]/(f_i(a_1, \ldots, a_{i-1}, X_i)).
\]

Now if $P(X) \in k[X_1, \ldots, X_n]$ is in the kernel of $\varphi$, we have $\varphi(P) =
P(a_1, \ldots, a_n) = 0$, so that $P(a_1, \ldots, a_{n-1}, X_n)$ is divisible by $f_n(a_1, \ldots,
a_{n-1}, X_n)$; dividing $P(X_1, \ldots, X_n)$ as a polynomial in $X_n$ by the
monic polynomial $f_n(X_1, \ldots, X_n)$ and letting $R_n(X_1, \ldots, X_n)$ be the remain-
der, we can write $P = Q_n f_n + R_n$, with $R_n(a_1, \ldots, a_{n-1}, X_n) = 0$. Similarly,
dividing $R_n(X_1, \ldots, X_n)$ as a polynomial in $X_{n-1}$ by $f_{n-1}(X_1, \ldots, X_{n-1})$ and
letting $R_{n-1}(X_1, \ldots, X_n)$ be the remainder, we get
\[
  R_n = Q_{n-1}f_{n-1} + R_{n-1},
\]
with
\[
  R_{n-1}(a_1, \ldots, a_{n-2}, X_{n-1}, X_n) = 0.
\]
Proceeding in the same way we get $P = \sum Q_i f_i + R$, with $R(X_1, \ldots, X_n) = 0$;
that is $R = 0$ and $P = \sum Q_i f_i$, so that $\Ker \varphi = (f_1, f_2, \ldots, f_n)$.
\end{proof}

The following theorem can be regarded as a converse of Theorem 1, (i).

\begin{theorem}\label{thm:5.2}
Let $k$ be a field and $A = k[\alpha_1, \ldots, \alpha_n]$ an integral domain, and
write $r = \trdeg_k A$ for the transcendence degree of $A$ (that is, of its field of
fractions) over $k$. Then if $r > 0$, $A$ is not a field.
\end{theorem}

\begin{proof}
Suppose that $\alpha_1, \ldots, \alpha_r$ is a transcendence basis for $A$ over $k$, and set
$K = k(\alpha_1, \ldots, \alpha_r)$. Then since $\alpha_{r+1}, \ldots, \alpha_n$ are algebraic over $K$, there exist
polynomials $f_i(X_{r+1}, \ldots, X_i) \in K[X_{r+1}, \ldots, X_i]$, monic of degree $d_i$ in $X_i$, such that
\[
  k[\alpha_{r+1}, \ldots, \alpha_n] \cong K[X_{r+1}, \ldots, X_n]/(f_{r+1}, \ldots, f_n)
\]
and
\[
  d_i = [K(\alpha_{r+1}, \ldots, \alpha_i) : K(\alpha_{r+1}, \ldots, \alpha_{i-1})] \quad \text{for } i > r.
\]
The coefficients of $f_i$ are in $K$, so that for suitable $0 \neq g \in k[\alpha_1, \ldots, \alpha_r]$ we
have $gf_i \in k[\alpha_1, \ldots, \alpha_r][X_{r+1}, \ldots, X_n]$. In other words, if we set $B =
k[\alpha_1, \ldots, \alpha_r, g^{-1}]$ then the $f_i$ are polynomials with coefficients in $B$. We
are now going to show that $A[g^{-1}] = B[\alpha_{r+1}, \ldots, \alpha_n]$ is a free module over
$B$ with basis $\{\prod_{i>r} \alpha_i^{e_i} \mid 0 \le e_i < d_i\}$. Every element of $B[\alpha_{r+1}, \ldots, \alpha_n]$ can
be written as $P(\alpha_{r+1}, \ldots, \alpha_n)$ for some $P \in B[X_{r+1}, \ldots, X_n]$; dividing $P$ as a
polynomial in $X_n$ by $f_n$, and replacing $P$ by the remainder, we can assume
that $P$ has degree at most $d_n - 1$ in $X_n$; then dividing $P$ as a polynomial
in $X_{n-1}$ by $f_{n-1}$, and replacing by the remainder, we can assume that $P$
has degree at most $d_{n-1} - 1$ in $X_{n-1}$. Proceeding in the same way, we
can assume that $P$ has degree at most $d_i - 1$ in $X_i$ for each $i$; in addition,
the elements $\{\alpha_i^e \mid 0 \le e < d_i\}$ are linearly independent over $K(\alpha_{r+1}, \ldots, \alpha_{i-1})$. Hence $A[g^{-1}]$ is a free $B$-module. However, $B$ is not a field; for
$k[a_1, \ldots, a_r]$ is a polynomial ring in $r$ variables over $k$, and hence it contains
infinitely many irreducible polynomials (the proof of this is exactly the
same as Euclid's proof that there exist infinitely many primes). Hence,
there is an irreducible polynomial $h \in k[a_1, \ldots, a_r]$ which does not divide
$g$, and then obviously $h^{-1} \notin k[a_1, \ldots, a_r, g^{-1}]$. Therefore $B$ contains an
ideal $I$ with $I \neq 0, B$, and since $A[g^{-1}]$ is a free module over $B$, $IA[g^{-1}]$
is a proper ideal of $A[g^{-1}]$. Thus $A[g^{-1}]$ is not a field. But if $A$ were a
field then we would have $A[g^{-1}] = A$, and hence $A$ is not a field.
\end{proof}

\begin{theorem}\label{thm:5.3}
Let $k$ be a field, and let $m$ be any maximal ideal of the
polynomial ring $k[X_1, \ldots, X_n]$; then the residue class field $k[X_1, \ldots, X_n]/m$
is algebraic over $k$. Hence $m$ can be generated by $n$ elements, and in
particular if $k$ is algebraically closed then $m$ is of the form $m =
(X_1 - a_1, \ldots, X_n - a_n)$ for $a_i \in k$.
\end{theorem}

\begin{proof}
Set $k[X_1, \ldots, X_n]/m = K$, and write $\alpha_i$ for the image of $X_i$ in $K$; then
$K = k[\alpha_1, \ldots, \alpha_n]$. By the previous theorem, since $K$ is a field it is algebraic
over $k$, and then by Theorem 1, (ii), $m$ is generated by $n$ elements. If $k$ is
algebraically closed then $k = K$, so that each $X_i$ is congruent modulo $m$ to
some $a_i \in k$; then $(X_1 - a_1, \ldots, X_n - a_n) \subset m$. On the other hand
$(X_1 - a_1, \ldots, X_n - a_n)$ is obviously a maximal ideal, so that equality must
hold.
\end{proof}

Let $k$ be a field and $\bar{k}$ its algebraic closure. Suppose that $\Phi \subset
k[X_1, \ldots, X_n]$ is a subset. An $n$-tuple $\alpha = (a_1, \ldots, a_n)$ of elements $a_i \in \bar{k}$ is an
\emph{algebraic zero}\index{algebraic zero} of $\Phi$ if it satisfies $f(\alpha) = 0$ for every $f(X) \in \Phi$.

\begin{theorem}[The Hilbert Nullstellensatz]\label{thm:5.4}
\begin{enumerate}[(i)]
  \item If $\Phi$ is a subset of $k[X_1, \ldots, X_n]$ which does not have any algebraic
        zeros then the ideal generated by $\Phi$ contains $1$.
  \item Given a subset $\Phi$ of $k[X_1, \ldots, X_n]$ and an element $f \in k[X_1, \ldots, X_n]$,
        suppose that $f$ vanishes at every algebraic zero of $\Phi$. Then some power of $f$
        belongs to the ideal generated by $\Phi$, that is there exist $\nu > 0$,
        $g_i \in k[X_1, \ldots, X_n]$ and $h_i \in \Phi$ such that $f^\nu = \sum g_i h_i$.
\end{enumerate}
\end{theorem}

\begin{proof}
(i) Let $I$ be the ideal generated by $\Phi$; if $1 \notin I$ then there exists a
maximal ideal $m$ containing $I$. By the previous theorem, $k[X_1, \ldots, X_n]/m$
is algebraic over $k$, so that it has a $k$-linear isomorphic embedding $\theta$ into
$\bar{k}$. If we set $\theta(X_i \bmod m) = a_i$ then for all $g(X) \in m$ we have $0 = \theta(g(X)) =
g(a)$, therefore $\alpha = (a_1, \ldots, a_n)$ is an algebraic zero of $m$, and
hence also of $\Phi$. This contradicts the hypothesis; hence, $1 \in I$.

(ii) Inside $k[X_1, \ldots, X_n, Y]$ we consider the set $\Phi \cup \{1 - Yf(X)\}$; then this
set has no algebraic zeros, so that by (i) it generates the ideal $(1)$. Therefore
there exists a relation of the form
\[
  1 = \sum p_i(X, Y)h_i(X) + Q(X, Y)(1 - Yf(X)),
\]
with $h_i(X) \in \Phi$. This is an identity in $X_1, \ldots, X_n$ and $Y$, so that it still holds if
we substitute $Y = f(X)^{-1}$. Hence we have
\[
  1 = \sum p_i(X, f^{-1})h_i(X),
\]
so that multiplying through by a suitable power of $f$ and clearing
denominators gives $f^\nu = \sum g_i(X)h_i(X)$, with $g_i \in k[X_1, \ldots, X_n]$ and
$h_i \in \Phi$.
\end{proof}

\begin{remark}
The above proof of (ii) is a classical idea due to Rabinowitch
[1]. In a modern form it can be given as follows: let $I \subset k[X_1, \ldots, X_n] = A$
be the ideal generated by $\Phi$; then in the localisation $A_f$ with respect to $f$ (see
\S\ref{sec:4}, Example 1), we have $IA_f = A_f$, so that a power of $f$ is in $I$.
\end{remark}

\begin{theorem}\label{thm:5.5}
Let $k$ be a field, $A$ a ring which is finitely generated over $k$, and
$I$ a proper ideal of $A$; then the radical of $I$ is the intersection of all maximal
ideals containing $I$, that is $\sqrt{I} = \bigcap_{I \subset m} m$.
\end{theorem}

\begin{proof}
Let $A = k[a_1, \ldots, a_n]$, so that $A$ is a quotient of $k[X_1, \ldots, X_n]$.
Considering the inverse image of $I$ in $k[X]$ reduces to the case $A = k[X]$,
and the assertion follows from Theorem 4, (ii).
\end{proof}

Compared to the result $\sqrt{I} = \bigcap_{I \subset P} P$ proved in \S\ref{sec:1}, the conclusion of
Theorem 5 is much stronger. It is equivalent to the condition on a ring
that every prime ideal $P$ should be expressible as an intersection of maximal
ideals. Rings for which this holds are called \emph{Hilbert rings}\index{Hilbert ring} or \emph{Jacobson rings}\index{Jacobson ring},
and they have been studied independently by O. Goldmann [1] and
W. Krull [7]. See also Kaplansky [K] and Bourbaki [B-5].

\begin{theorem}\label{thm:5.6}
Let $k$ be a field and $A$ an integral domain which is finitely
generated over $k$; then
\[
  \dim A = \trdeg_k A.
\]
\end{theorem}

\begin{proof}
Let $A = k[X_1, \ldots, X_n]/P$, and set $r = \trdeg_k A$. To prove that
$r \ge \dim A$ it is enough to show that if $P$ and $Q$ are prime ideals of
$k[X] = k[X_1, \ldots, X_n]$ with $Q \supset P$ and $Q \neq P$ then
\[
  \trdeg_k k[X]/Q < \trdeg_k k[X]/P.
\]
The $k$-algebra homomorphism $k[X]/P \to k[X]/Q$ is onto, so that $\trdeg_k
k[X]/Q \le \trdeg_k k[X]/P$ is obvious. Suppose that equality holds. Let
$k[X]/P = k[\alpha_1, \ldots, \alpha_n]$ and $k[X]/Q = k[\beta_1, \ldots, \beta_n]$; we can assume that
$\beta_1, \ldots, \beta_r$ is a transcendence basis for $k(\beta)$ over $k$. Then $\alpha_1, \ldots, \alpha_r$ are also
algebraically independent over $k$, so that they form a transcendence basis
for $k(\alpha)$ over $k$. Now set $S = k[X_1, \ldots, X_n] - \{0\}$; $S$ is a multiplicative set in
$k[X]$ with $P \cap S = \emptyset$ and $Q \cap S = \emptyset$. Setting $R = k[X_1, \ldots, X_n]$ and
$K = k(X_1, \ldots, X_n)$, we have $R_S = K[X_{r+1}, \ldots, X_n]$, and
\[
  R_S/PR_S \cong k(\alpha_1, \ldots, \alpha_r)[\alpha_{r+1}, \ldots, \alpha_n],
\]
so that $R_S/PR_S$ is algebraic over $K = k(X_1, \ldots, X_n) \cong k(\alpha_1, \ldots, \alpha_r)$, and
therefore by Theorem 1, $PR_S$ is a maximal ideal of $R_S$; but this contradicts
the assumptions $P \subset Q$ with $P \neq Q$ and $Q \cap S = \emptyset$.

Now let us prove that $r \le \dim A$ by induction on $r$. If $r = 0$ then, by
Theorem 1, $A$ is a field, so $\dim A = 0$ and the assertion holds. Now let $r > 0$,
and suppose that $A = k[\alpha_1, \ldots, \alpha_n]$ with $\alpha_1$ transcendental over $k$; setting
$S = k[X_1] - \{0\}$ and $R = k[X_1, \ldots, X_n]$ we get
\[
  R_S = k(X_1)[X_2, \ldots, X_n] \quad \text{and} \quad R_S/PR_S \cong k(X_1)[\alpha_2, \ldots, \alpha_n].
\]
Hence $R_S/PR_S$ has transcendence degree $r - 1$ over $k(X_1)$, so that by
induction $\dim R_S/PR_S \ge r - 1$. Thus there exists a strictly increasing chain
$PR_S = Q_0 \subsetneq Q_1 \subsetneq \cdots \subsetneq Q_{r-1}$ of prime ideals of $R_S$. If we set $P_i = Q_i \cap R$ then
$P_i$ is a prime ideal of $R$ disjoint from $S$; in particular, the residue class of
$X_1$ in $R/P_{r-1}$ is not algebraic over $k$, and so $\trdeg_k R/P_{r-1} \ge 1$. Then
$P_{r-1}$ is not a maximal ideal of $R$ by Theorem 3, and therefore $R$ has a
maximal ideal $P_r$ strictly bigger than $P_{r-1}$. Hence $\dim A = \coht P \ge r$.
\end{proof}

\begin{corollary*}
If $k$ is a field then $\dim k[X_1, \ldots, X_n] = n$.
\end{corollary*}

We now turn to a different topic, the theorem of Forster and Swan on the
number of generators of a module. Let $A$ be a ring and $M$ a finite $A$-module;
for $p \in \Spec A$, write $\kappa(p)$ for the residue field of $A_p$, and let $\mu(p, M)$ denote the
dimension over $\kappa(p)$ of the vector space $M \otimes \kappa(p) = M_p/pM_p$ (in the usual
sense of linear algebra). This is the cardinality of a minimal basis of the
$A_p$-module $M_p$. Hence, if $p \supset p'$ then $\mu(p, M) \ge \mu(p', M)$.

In 1964 the young function-theorist O. Forster surprised the experts in
algebra by proving the following theorem [1].

\begin{theorem}\label{thm:5.7}
Let $A$ be a Noetherian ring and $M$ a finite $A$-module. Set
\[
  b(M) = \sup\{\mu(p, M) + \coht p \mid p \in \Supp M\};
\]
then $M$ can be generated by at most $b(M)$ elements.
\end{theorem}

This theorem is a very important link between the number of local and
global generators. However, there was room for improvement in the bound
for the number of generators, and in no time R. Swan obtained a better
bound. We will prove Swan's bound. For this we need the concept of \emph{j-Spec
A}\index{j-Spec of a ring} introduced by Swan. This is a space having the same irreducible closed
subsets as m-Spec $A$, but has the advantage of having a `generic point', not
present in m-Spec $A$, for every irreducible closed subset.

A prime ideal which can be expressed as an intersection of maximal
ideals is called a \emph{$j$-prime ideal}\index{j-prime ideal}, and we write j-Spec $A$ for the set of all
$j$-prime ideals. We consider j-Spec $A$ also with its topology as a subspace
of $\Spec A$. Set $M = \text{m-Spec } A$ and $J = \text{j-Spec } A$. If $F$ is a closed subset of
$J$ then there is an ideal $I$ of $A$ such that $F = V(I) \cap J$. One sees easily that
a prime ideal $P$ belongs to $F$ if and only if $P$ can be expressed as an
intersection of elements of $F \cap M = V(I) \cap M$. Hence $F$ is determined by
$F \cap M$, so that there is a natural one-to-one correspondence between
closed subsets of $J$ and of $M$. It follows that if $M$ is Noetherian so is $J$,
and they both have the same combinatorial dimension. Now let $B$ be an
irreducible closed subset of $J$, and let $P$ be the intersection of all the
elements of $B$. If $B = V(I) \cap J$ then $I \subset P$ and we can also write
$B = V(P) \cap J$. If $P$ is not a prime ideal then there exist $f, g \in A$ such that
$f \notin P$, $g \notin P$ and $fg \in P$; but then
\[
  B = (V(P + fA) \cap J) \cup (V(P + gA) \cap J),
\]
and by definition of $P$ there is a $Q \in B$ not containing $f$ and a $Q' \in B$ not
containing $g$, which implies that $B$ is reducible, a contradiction. Therefore
$P$ is a prime ideal. Hence $P \in B$ and $B = V(P) \cap J$. This $P$ is called the
\emph{generic point}\index{generic point} of $B$. Conversely if $P$ is any element of $J$ then $V(P) \cap J$ is
an irreducible closed subset of $J$, and is the closure in $J$ of $\{P\}$. We will
write \emph{$j$-dim $p$} for the combinatorial dimension of $V(p) \cap J$.

For a finite $A$-module $M$ and $p \in J$ we set
\[
  b(p, M) =
  \begin{cases}
    0 & \text{if } M_p = 0 \\
    \text{$j$-dim } p + \mu(p, M) & \text{if } M_p \neq 0.
  \end{cases}
\]

\begin{theorem}[Swan~\cite{Sw1}]\label{thm:5.8}
Let $A$ be a ring, and suppose that m-Spec $A$ is
a Noetherian space. Let $M$ be a finite $A$-module. If
\[
  \sup\{b(p, M) \mid p \in \text{j-Spec } A\} = r < \infty
\]
then $M$ is generated by at most $r$ elements.
\end{theorem}

\begin{proof}
\textbf{Step 1.} For $p \in \Spec A$ and $x \in M$, we will say that $x$ is \emph{basic} at $p$ if $x$
has non-zero image in $M \otimes \kappa(p)$. It is easy to see that this condition is
equivalent to $\mu(p, M/Ax) = \mu(p, M) - 1$.

\begin{lemma*}
Let $M$ be a finite $A$-module and $p_1, \ldots, p_n \in \Supp(M)$. Then there
exists $x \in M$ which is basic at each of $p_1, \ldots, p_n$.
\end{lemma*}

\begin{proof}[Proof of the lemma]
By reordering $p_1, \ldots, p_n$ we assume that $p_i$ is maximal among
$\{p_i, p_{i+1}, \ldots, p_n\}$ for each $i$. By induction on $n$ suppose that $x' \in M$ is basic at
$p_1, \ldots, p_{n-1}$. If $x'$ is basic at $p_n$, then we can take $x = x'$. Suppose then that $x'$
is not basic at $p_n$. By assumption $M_{p_n} \neq 0$ so that we can choose some $y \in M$
which is basic at $p_n$. We have $p_1 \cdots p_{n-1} \not\subset p_n$, so that if we take an element
$a \in p_1 \cdots p_{n-1}$ not belonging to $p_n$, and set $x = x' + ay$, this $x$ satisfies our
requirements. This proves the lemma.
\end{proof}

\textbf{Step 2.} Setting $\sup\{b(p, M) \mid p \in \text{j-Spec } A\} = r$, we now show that there are
just a finite number of primes $p$ such that $b(p, M) = r$. Indeed, for
$n = 1, 2, \ldots$, the subset $X_n = \{p \in \text{j-Spec } A \mid \mu(p, M) \ge n\}$ is closed in j-Spec $A$ by
Theorem~\ref{thm:4.10}; it has a finite number of irreducible components (by
Ex.~\ref{ex:4.11-ex}), and we let $p_{ni}$ (for $1 \le i \le v_n$) be their generic points. If $M$ is
generated by $s$ elements then $X_n = \emptyset$ for $n > s$, so that the set $\{p_{ni}\}_{n,i}$ is
finite. Let us prove that if $b(p, M) = r$ then $p \in \{p_{ni}\}_{n,i}$. Suppose
$\mu(p, M) = n$; then $p \in X_n$, so that by construction $p \supset p_{ni}$ for some $i$. But
if $p \neq p_{ni}$ then j-dim $p$ < j-dim $p_{ni}$, and since $\mu(p, M) = n = \mu(p_{ni}, M)$ we have
$b(p, M) < b(p_{ni}, M)$, which is a contradiction. Hence $p = p_{ni}$.

\textbf{Step 3.} Let us choose an element $x \in M$ which is basic for each of the
finitely many primes $p$ with $b(p, M) = r$, and set $\bar{M} = M/Ax$; then clearly
$b(p, \bar{M}) \le r - 1$ for every $p \in \text{j-Spec } A$. Hence by induction $\bar{M}$ is generated by
$r - 1$ elements, and therefore $M$ by $r$ elements.
\end{proof}

Swan's paper contains a proof of the following generalisation to non-
commutative rings: let $A$ be a commutative ring, $\Lambda$ a possibly non-
commutative $A$-algebra and $M$ a finite left $\Lambda$-module. Suppose that
m-Spec $A$ is Noetherian, and that for every maximal ideal $p$ of $A$ the
$\Lambda_p$-module $M_p$ is generated by at most $r$ elements; then $M$ is generated as a
$\Lambda$-module by at most $r + d$ elements, where $d$ is the combinatorial
dimension of m-Spec $A$.

The Forster--Swan theorem is a statement that local properties imply
global ones; remarkable results in this direction have been obtained by
Mohan Kumar [2] (see also Cowsik--Nori [1] and Eisenbud--Evans [1],
[2]). The number of generators of ideals in local rings is the subject of
a nice book by J. Sally [Sa].

% ============ 习题 ============
\section*{Exercises to \S\ref{sec:5}}
\emph{Prove the following propositions.}

\begin{exercise}\label{ex:5.1-ex}
Let $k$ be a field, $R = k[X_1, \ldots, X_n]$ and let $P \in \Spec R$; then $\htop P +
\coht P = n$.
\end{exercise}

\begin{exercise}\label{ex:5.2-ex}
A zero-dimensional Noetherian ring is Artinian (the converse to Example
2 above).
\end{exercise}

% ============ §6 Associated primes and primary decomposition (原书 pp.37-42) ============
\section{Associated primes and primary decomposition}\label{sec:6}\index{associated primes, section 6}\index{primary decomposition, section 6}

Most readers will presumably have come across primary decom-
position of ideals in Noetherian rings. This was the first big theorem
obtained by Emmy Noether in her abstract treatment of commutative
rings. Nowadays, as exemplified by Bourbaki [B4], the notion of associated
prime is considered more important than primary decomposition itself.

Let $A$ be a ring and $M$ an $A$-module. A prime ideal $P$ of $A$ is called an
\emph{associated prime ideal}\index{associated prime ideal} of $M$ if $P$ is the annihilator $\ann(x)$ of some $x \in M$. The
set of associated primes of $M$ is written $\Ass(M)$ or $\Ass_A(M)$. For $I$ an ideal
of $A$, the associated primes of the $A$-module $A/I$ are referred to as the
\emph{prime divisors}\index{prime divisor} of $I$. We say that $a \in A$ is a \emph{zero-divisor} for $M$ if there is a non-zero
$x \in M$ such that $ax = 0$, and otherwise that $a$ is \emph{$M$-regular}.\index{regular element}

\begin{theorem}\label{thm:6.1}
Let $A$ be a Noetherian ring and $M$ a non-zero $A$-module.
\begin{enumerate}[(i)]
  \item Every maximal element of the family of ideals $\mathcal{F} = \{\ann(x) \mid 0 \neq x \in M\}$
        is an associated prime of $M$, and in particular $\Ass(M) \neq \emptyset$.
  \item The set of zero-divisors for $M$ is the union of all the associated primes
        of $M$.
\end{enumerate}
\end{theorem}

\begin{proof}
(i) We have to prove that if $\ann(x)$ is a maximal element of $\mathcal{F}$ then it is
prime: if $a, b \in A$ are such that $abx = 0$ but $bx \neq 0$ then by maximality
$\ann(bx) = \ann(x)$; hence, $ax = 0$.

(ii) If $ax = 0$ for some $x \neq 0$ then $a \in \ann(x) \in \mathcal{F}$, and by (i) there is an
associated prime of $M$ containing $\ann(x)$.
\end{proof}

\begin{theorem}\label{thm:6.2}
Let $S \subset A$ be a multiplicative set, and $N$ an $A_S$-module.
Viewing $\Spec(A_S)$ as a subset of $\Spec A$, we have $\Ass_A(N) = \Ass_{A_S}(N)$. If $A$
is Noetherian then for an $A$-module $M$ we have $\Ass(M_S) =
\Ass(M) \cap \Spec(A_S)$.
\end{theorem}

\begin{proof}
For $x \in N$ we have $\ann_A(x) = \ann_{A_S}(x) \cap A$, so that if $P \in \Ass_{A_S}(N)$
then $P \cap A \in \Ass_A(N)$. Conversely if $p \in \Ass_A(N)$ and we choose $x \in N$
such that $p = \ann_A(x)$ then $x \neq 0$, and hence, $p \cap S = \emptyset$ and $pA_S$ is a prime
ideal of $A_S$ with $pA_S = \ann_{A_S}(x)$. For the second part, if $p \in \Ass(M) \cap
\Spec(A_S)$ then $p \cap S = \emptyset$, and $p = \ann_A(x)$ for some $x \in M$; if $(a/s)x = 0$ in
$M_S$ then there is a $t \in S$ such that $tax = 0$ in $M$, and $t \notin p$, $ta \in p$ gives
$a \in p$, so that $\ann_{A_S}(x) = pA_S$ and $pA_S \in \Ass(M_S)$. Conversely, if $P \in
\Ass(M_S)$ then without loss of generality we have $P = \ann_{A_S}(x)$ with $x \in M$.
Setting $p = P \cap A$ we have $P = pA_S$. Now $p$ is finitely generated since $A$ is
Noetherian, and it follows that there exists some $t \in S$ such that $p =
\ann_A(tx)$. Therefore $p \in \Ass_A(M)$.
\end{proof}

\begin{corollary*}
For a Noetherian ring $A$, an $A$-module $M$ and a prime ideal $P$ of
$A$ we have
\[
  P \in \Ass_A(M) \iff PA_P \in \Ass_{A_P}(M_P).
\]
\end{corollary*}

\begin{theorem}\label{thm:6.3}
Let $A$ be a ring and $0 \to M' \to M \to M'' \to 0$ an exact
sequence of $A$-modules; then
\[
  \Ass(M) \subset \Ass(M') \cup \Ass(M'').
\]
\end{theorem}

\begin{proof}
If $P \in \Ass(M)$ then $M$ contains a submodule $N$ isomorphic to $A/P$.
Since $P$ is prime, for any non-zero element $x$ of $N$ we have $\ann(x) = P$.
Therefore if $N \cap M' \neq 0$ we have $P \in \Ass(M')$. If $N \cap M' = 0$ then the image
of $N$ in $M''$ is also isomorphic to $A/P$, so that $P \in \Ass(M'')$.
\end{proof}

\begin{theorem}\label{thm:6.4}
Let $A$ be a Noetherian ring and $M \neq 0$ a finite $A$-module.
Then there exists a chain $0 = M_0 \subset M_1 \subset \cdots \subset M_r = M$ of submodules of $M$
such that for each $i$ we have $M_i/M_{i-1} \cong A/P_i$ with $P_i \in \Spec A$.
\end{theorem}

\begin{proof}
Choose any $P_1 \in \Ass(M)$; then there exists a submodule $M_1$ of $M$
with $M_1 \cong A/P_1$. If $M_1 \neq M$ and we choose any $P_2 \in \Ass(M/M_1)$ then
there exists $M_2 \subset M$ such that $M_2/M_1 \cong A/P_2$. Continuing in the same way and
using the ascending chain condition, we eventually arrive at $M_r = M$.
\end{proof}

\begin{theorem}\label{thm:6.5}
Let $A$ be a Noetherian ring and $M$ a finite $A$-module.
\begin{enumerate}[(i)]
  \item $\Ass(M)$ is a finite set.
  \item $\Ass(M) \subset \Supp(M)$.
  \item The set of minimal elements of $\Ass(M)$ and of $\Supp(M)$ coincide.
\end{enumerate}
\end{theorem}

\begin{proof}
(i) follows from the previous two theorems; we need only note
that $\Ass(A/P) = \{P\}$. For (ii), if $0 \to A/P \to M$ is exact then so is $0 \to
A_P/PA_P \to M_P$, and therefore $M_P \neq 0$. For (iii) it is enough to show that
if $P$ is a minimal element of $\Supp(M)$ then $P \in \Ass(M)$. We have $M_P \neq 0$
so that by Theorem 2 and (ii),
\[
  \emptyset \neq \Ass(M_P) = \Ass(M) \cap \Spec(A_P) \subset \Supp(M) \cap \Spec(A_P)
             = \{P\}.
\]
Therefore we must have $P \in \Ass(M)$.
\end{proof}

Let $A$ be a Noetherian ring and $M$ a finite $A$-module. Let $P_1, \ldots, P_r$ be the
minimal elements of $\Supp(M)$; then $\Supp(M) = V(P_1) \cup \cdots \cup V(P_r)$, and the
$V(P_i)$ are the irreducible components of the closed set $\Supp(M)$ (see
Ex.~\ref{ex:4.11-ex}). The prime ideals $P_1, \ldots, P_r$ are called the \emph{isolated associated primes}\index{isolated associated prime}
of $M$, and the remaining associated primes of $M$ are called \emph{embedded}
\emph{primes}.\index{embedded prime} If $I$ is an ideal of $A$ then $\Supp_A(A/I)$ is the set of prime ideals
containing $I$, and the \emph{minimal prime divisors} of $I$ (that is the minimal
associated primes of the $A$-module $A/I$) are precisely the minimal prime
ideals containing $I$. We have seen in Ex.~\ref{ex:4.12-ex} that there are only a finite
number of such primes, and Theorem 5 now gives a new proof of this.
(For examples of embedded primes see Ex.~\ref{ex:6.6-ex} and Ex.~\ref{ex:8.9-ex}.)

\begin{definition}
Let $A$ be a ring, $M$ an $A$-module and $N \subset M$ a submodule. We
say that $N$ is a \emph{primary submodule}\index{primary submodule} of $M$ if the following condition holds for
all $a \in A$ and $x \in M$:
\[
  x \notin N \text{ and } ax \in N \Rightarrow a^\nu M \subset N \text{ for some } \nu.
\]
This definition in fact only depends on the quotient module $M/N$. It can be
restated as
\[
  \text{if } a \in A \text{ is a zero-divisor for } M/N \text{ then } a \in \sqrt{\ann(M/N)}.
\]
\end{definition}

A primary ideal is just a primary submodule of the $A$-module $A$. One
might wonder about trying to set up a notion of prime submodule
generalising prime ideal, but this does not turn out to be useful.

\begin{theorem}\label{thm:6.6}
Let $A$ be a Noetherian ring and $M$ a finite $A$-module. Then a
submodule $N \subset M$ is primary if and only if $\Ass(M/N)$ consists of one
element only. In this case, if $\Ass(M/N) = \{P\}$ and $\ann(M/N) = I$ then $I$ is
primary and $\sqrt{I} = P$.
\end{theorem}

\begin{proof}
If $\Ass(M/N) = \{P\}$ then by the previous theorem $\Supp(M/N) =
V(P)$, so that $P = \sqrt{\ann(M/N)}$. Now if $a \in A$ is a zero-divisor for $M/N$ it
follows from Theorem 1 that $a \in P$, so that $a \in \sqrt{\ann(M/N)}$; hence, $N$ is a
primary submodule of $M$. Conversely, if $N$ is a primary submodule and
$P \in \Ass(M/N)$ then every $a \in P$ is a zero-divisor for $M/N$, so that by
assumption $a \in \sqrt{I}$, where $I = \ann(M/N)$. Hence $P \subset \sqrt{I}$, but from the
definition of associated prime we obviously have $I \subset P$, and hence $\sqrt{I} \subset P$,
so that $P = \sqrt{I}$. Thus $\Ass(M/N)$ has just one element $\sqrt{I}$. We prove that in
this case $I$ is a primary ideal: let $a, b \in A$ with $b \notin I$; if $ab \in I$ then $ab(M/N) = 0$,
but $b(M/N) \neq 0$, so that $a$ is a zero-divisor for $M/N$, and therefore
$a \in P = \sqrt{I}$.
\end{proof}

\begin{definition}
If $\Ass(M/N) = \{P\}$ we say that $N \subset M$ is a \emph{$P$-primary} sub-
module, or a \emph{primary submodule belonging to $P$}.
\end{definition}

\begin{theorem}\label{thm:6.7}
If $N$ and $N'$ are $P$-primary submodules of $M$ then so is $N \cap N'$.
\end{theorem}

\begin{proof}
We can embed $M/(N \cap N')$ as a submodule of $(M/N) \oplus (M/N')$, so
that
\[
  \Ass(M/(N \cap N')) \subset \Ass(M/N) \cup \Ass(M/N') = \{P\}.
\]
\end{proof}

If $N \subset M$ is a submodule, we say that $N$ is \emph{reducible}\index{reducible submodule} if it can be written as
an intersection $N = N_1 \cap N_2$ of two submodules $N_1, N_2$ with $N_i \neq N$, and
otherwise that $N$ is \emph{irreducible}; note that this has nothing to do with the
notion of irreducible modules in representation theory (= no submodules
other than $0$ and $M$), which is a condition on $M$ only.

If $M$ is a Noetherian module then any submodule $N$ of $M$ can be written
as a finite intersection of irreducible submodules. Proof: let $\mathcal{F}$ be the set
of submodules $N \subset M$ having no such expression. If $\mathcal{F} \neq \emptyset$ then it has a
maximal element $N_0$. Then $N_0$ is reducible, so that $N_0 = N_1 \cap N_2$, and
$N_i \notin \mathcal{F}$. Now each of the $N_i$ is an intersection of a finite number of
irreducible submodules, and hence so is $N_0$. This is a contradiction.

\begin{remark}
The representation as an intersection of irreducible submodules is
in general not unique. For example, if $A$ is a field and $M$ an $n$-dimensional
vector space over $A$ then the irreducible submodules of $M$ are just its
$(n - 1)$-dimensional subspaces. An $(n - 2)$-dimensional subspace can be
written in lots of ways as an intersection of $(n - 1)$-dimensional subspaces.
\end{remark}

In general we say that an expression of a set $N$ as an intersection
$N = N_1 \cap \cdots \cap N_r$ is \emph{irredundant}\index{irredundant decomposition} if we cannot omit any $N_i$, that is if
$N \neq N_1 \cap \cdots \cap N_{i-1} \cap N_{i+1} \cap \cdots \cap N_r$. If $M$ is an $A$-module, we call an
expression $N = N_1 \cap \cdots \cap N_r$ of a submodule $N$ as an intersection of a finite
number of submodules $N_i \subset M$ a \emph{decomposition} of $N$; if each of the $N_i$ is
irreducible we speak of an \emph{irreducible decomposition}, if primary of a \emph{primary decomposition}.
Let $N = N_1 \cap \cdots \cap N_r$ be an irredundant primary decompo-
sition with $\Ass(M/N_i) = \{P_i\}$; if $P_i = P_j$ then $N_i \cap N_j$ is again primary, so
that grouping together all of the $N_i$ belonging to the same prime ideal we get
a primary decomposition such that $P_i \neq P_j$ for $i \neq j$. A decomposition with
this property will be called a \emph{shortest primary decomposition}, and the $N_i$
appearing in it the \emph{primary components} of $N$; if $N_i$ belongs to a prime $P$ we
sometimes say that $N_i$ is the \emph{$P$-primary component} of $N$.

\begin{theorem}\label{thm:6.8}
Let $A$ be a Noetherian ring and $M$ a finite $A$-module.
\begin{enumerate}[(i)]
  \item An irreducible submodule of $M$ is a primary submodule.
  \item If
        \[
          N = N_1 \cap \cdots \cap N_l \quad \text{with } \Ass(M/N_i) = \{P_i\}
        \]
        is an irredundant primary decomposition of a proper submodule $N \subset M$
        then $\Ass(M/N) = \{P_1, \ldots, P_l\}$.
  \item Every proper submodule $N$ of $M$ has a primary decomposition. If $N$ is
        a proper submodule of $M$ and $P$ is a minimal associated prime of $M/N$ then
        the $P$-primary component of $N$ is $\varphi_P^{-1}(N_P)$, where $\varphi_P\colon M \to M_P$ is the
        canonical map, and therefore it is uniquely determined by $M$, $N$ and $P$.
\end{enumerate}
\end{theorem}

\begin{proof}
(i) It is enough to prove that a submodule $N \subset M$ which is not
primary is reducible: replacing $M$ by $M/N$ we can assume that $N = 0$. By
Theorem 6, $\Ass(M)$ has at least two elements $P_1$ and $P_2$. Then $M$ contains
submodules $K_i$ isomorphic to $A/P_i$ for $i = 1, 2$. Now since $\ann(x) = P_i$ for
any non-zero $x \in K_i$, we must have $K_1 \cap K_2 = 0$, and hence $0$ is reducible.

(ii) We can again assume that $N = 0$. If $0 = N_1 \cap \cdots \cap N_l$ then $M$ is
isomorphic to a submodule of $M/N_1 \oplus \cdots \oplus M/N_l$, so that
\[
  \Ass(M) \subset \bigcup_i \Ass(M/N_i) = \{P_1, \ldots, P_l\}.
\]
On the other hand $N_2 \cap \cdots \cap N_l \neq 0$, and taking $0 \neq x \in N_2 \cap \cdots \cap N_l$, we
have $\ann(x) = 0:x \subset N_1:x$. But $N_1:M$ is a primary ideal belonging to $P_1$, so
that $P_1^\nu M \subset N_1$ for some $\nu > 0$. Therefore $P_1^\nu x = 0$; hence there exists $i \ge 0$
such that $P_1^i x \neq 0$ but $P_1^{i+1} x = 0$, and choosing $0 \neq y \in P_1^i x$ we have
$P_1y = 0$. However, since $y \in N_2 \cap \cdots \cap N_l$ it follows that $y \notin N_1$, and by
the definition of primary submodule $\ann(y) \subset P_1$, so that $P_1 = \ann(y)$ and
$P_1 \in \Ass(M)$. The same works for the other $P_i$, and this proves that
$\{P_1, \ldots, P_l\} \subset \Ass(M)$.

(iii) We have already seen that a proper submodule has an irreducible
decomposition, so that by (i) it has a primary decomposition. Suppose that
$N = N_1 \cap \cdots \cap N_l$ is a shortest primary decomposition, and that $N_1$ is the
$P$-primary component with $P = P_1$. By Ex.~\ref{ex:4.8-ex} we know that $N_P =
(N_1)_P \cap \cdots \cap (N_l)_P$, and for $i > 1$ a power of $P_i$ is contained in $\ann(M/N_i)$;
then since $P_i \neq P_1$ we have $(M/N_i)_P = 0$, and therefore $(N_i)_P = M_P$. Thus
$N_P = (N_1)_P$ and hence $\varphi_P^{-1}(N_P) = \varphi_P^{-1}((N_1)_P)$; it is easy to check that
the right-hand side is $N_1$.
\end{proof}

\begin{remark}
The uniqueness of the $P$-primary component $N_1$ proved in (iii) for
minimal primes $P$ does not hold in general; see Ex.~\ref{ex:6.6-ex}.
\end{remark}

% ============ 习题 ============
\section*{Exercises to \S\ref{sec:6}}

\begin{exercise}\label{ex:6.1-ex}
Find $\Ass(M)$ for the $\Z$-module $M = \Z \oplus (\Z/3\Z)$.
\end{exercise}

\begin{exercise}\label{ex:6.2-ex}
If $M$ is a finite module over a Noetherian ring $A$, and $M_1, M_2$ are
submodules of $M$ with $M = M_1 + M_2$ then can we say that $\Ass(M) =
\Ass(M_1) \cup \Ass(M_2)$?
\end{exercise}

\begin{exercise}\label{ex:6.3-ex}
Let $A$ be a Noetherian ring and let $x \in A$ be an element which is neither a
unit nor a zero-divisor; prove that the ideals $xA$ and $x^nA$ for $n = 1, 2, \ldots$
have the same prime divisors:
\[
  \Ass_A(A/xA) = \Ass_A(A/x^nA).
\]
\end{exercise}

\begin{exercise}\label{ex:6.4-ex}
Let $I$ and $J$ be ideals of a Noetherian ring $A$. Prove that if $JA_P \subset IA_P$ for
every $P \in \Ass_A(A/I)$ then $J \subset I$.
\end{exercise}

\begin{exercise}\label{ex:6.5-ex}
Prove that the total ring of fractions of a reduced Noetherian ring $A$ is a
direct product of fields.
\end{exercise}

\begin{exercise}\label{ex:6.6-ex}
(Taken from [Nor 1], p. 30.) Let $k$ be a field. Show that in $k[X, Y]$ we have
\[
  (X^2, XY) = (X) \cap (X^2, Y) = (X) \cap (X^2, XY, Y^2).
\]
\end{exercise}

\begin{exercise}\label{ex:6.7-ex}
Let $f\colon A \to B$ be a homomorphism of Noetherian rings, and $M$ a finite $B$-
module. Write ${}^a f\colon \Spec B \to \Spec A$ as in \S\ref{sec:4}. Prove that ${}^a f(\Ass_B(M)) =
\Ass_A(M)$. (Consequently, $\Ass_A(M)$ is a finite set for such $M$.)
\end{exercise}

% ============ Appendix to §6. Secondary representations of a module (原书 pp.42-44) ============
\section*{Appendix to \S\ref{sec:6}. Secondary representations of a module}
\phantomsection
\addcontentsline{toc}{section}{Appendix to \S\ref{sec:6}. Secondary representations of a module}

I.G. Macdonald~\cite{MacD1} has developed the theory of \emph{attached prime ideals}\index{attached prime ideal}
and \emph{secondary representations}\index{secondary representation} of a module, which is in a certain sense
dual to the theory of associated prime ideals and primary decompositions.
This theory was successfully applied to the theory of local cohomology
by him and R.Y. Sharp (Macdonald \& Sharp~\cite{MS1}, Sharp~\cite{Shar7}).

Let $A$ be a commutative ring. An $A$-module $M$ is said to be \emph{secondary}\index{secondary module}
if $M \neq 0$ and, for each $a \in A$, the endomorphism $\varphi_a\colon M \to M$ defined
by $\varphi_a(m) = am$ (for $m \in M$) is either surjective or nilpotent. If $M$ is secondary,
then $P = \sqrt{\ann M}$ is a prime ideal, and $M$ is said to be \emph{$P$-secondary}. Any
non-zero quotient of a $P$-secondary module is $P$-secondary.

\begin{example}\label{ex:6.5}
If $A$ is an integral domain, its quotient field $K$ is a $(0)$-secondary
$A$-module.
\end{example}

\begin{example}\label{ex:6.6}
Let $W = \Z[p^{-1}]$, where $p$ is a prime number, and consider the
Artinian $\Z$-module $W/\Z$ (see \S\ref{sec:3}). This is also a $(0)$-secondary $\Z$-module.
\end{example}

\begin{example}\label{ex:6.7}
If $A$ is a local ring with maximal ideal $P$ and if every element of
$p$ is nilpotent, then $A$ itself is a $P$-secondary $A$-module.
\end{example}

\begin{example}\label{ex:6.8}
If $P$ is a maximal ideal of $A$, then $A/P^n$ is a $P$-secondary $A$-
module for every $n > 0$.
\end{example}

A \emph{secondary representation} of an $A$-module $M$ is an expression of $M$ as a
finite sum of secondary submodules:
\[
  (*) \quad M = N_1 + \cdots + N_r.
\]
The representation is \emph{minimal}\index{minimal secondary representation} if (1) the prime ideals $P_i := \sqrt{\ann N_i}$ are all
distinct, and (2) none of the $N_i$ is redundant. It is easy to see that the sum of
two $P$-secondary submodules is again $P$-secondary, hence if $M$ has a
secondary representation then it has a minimal one.

A prime ideal $P$ is called an \emph{attached prime ideal} of $M$ if $M$ has a $P$-
secondary quotient. The set of the attached prime ideals of $M$ is denoted by
$\Att(M)$.

\begin{theorem}\label{thm:6.9}
If $(*)$ is a minimal secondary representation of $M$ and $P_i =
\sqrt{\ann N_i}$, then $\Att(M) = \{P_1, \ldots, P_n\}$.
\end{theorem}

\begin{proof}
Since $M/(N_1 + \cdots + N_{i-1} + N_{i+1} + \cdots + N_n)$ is a non-zero quotient
of $N_i$, it is a $P_i$-secondary module. Thus $\{P_1, \ldots, P_n\} \subset \Att(M)$. Conver-
sely, let $P \in \Att(M)$ and let $W$ be a $P$-secondary quotient of $M$. Then $W =
W_1 + \cdots + W_n$, where $W_i$ is the image of $N_i$ in $W$. From this we obtain a
minimal secondary representation $W = W'_1 + \cdots + W'_m$, and then
$\Att(W) \supset \{P'_1, \ldots, P'_m\}$. On the other hand $\Att(W) = \{P\}$ since $W$ is
$P$-secondary. Therefore $P = P_i$ for some $i$.
\end{proof}

\begin{theorem}\label{thm:6.10}
If $0 \to M' \to M \to M'' \to 0$ is an exact sequence of
$A$-modules, then $\Att(M'') \subset \Att(M) \subset \Att(M') \cup \Att(M'')$.
\end{theorem}

\begin{proof}
The first inclusion is trivial from the definition. For the second, let
$P \in \Att(M)$ and let $N$ be a submodule such that $M/N$ is $P$-secondary. If
$M' + N = M$ then $M/N$ is a non-trivial quotient of $M'$, hence $P \in \Att(M')$.
If $M' + N \neq M$ then $M/(M' + N)$ is a non-trivial quotient of $M''$ as well
as of $M/N$, hence $M''$ has a $P$-secondary quotient and $P \in \Att(M'')$.
\end{proof}

An $A$-module $M$ is said to be \emph{sum-irreducible}\index{sum-irreducible module} if it is neither zero nor the
sum of two proper submodules.

\begin{lemma*}\label{lem:6.1}
If $M$ is Artinian and sum-irreducible, then it is secondary.
\end{lemma*}

\begin{proof}
Suppose $M$ is not secondary. Then there is $a \in A$ such that $M \neq aM$
and $a^nM \neq 0$ for all $n > 0$. Since $M$ is Artinian, we have $a^nM = a^{n+1}M$ for
some $n$. Set $K = \{x \in M \mid a^nx = 0\}$. Then it is immediate that $M = K + aM$,
and so $M$ is not sum-irreducible.
\end{proof}

\begin{theorem}\label{thm:6.11}
If $M$ is Artinian, then it has a secondary representation.
\end{theorem}

\begin{proof}
Similar to the proof of Theorem~\ref{thm:6.8}, (iii).
\end{proof}

The class of modules which have secondary representations is larger
than that of Artinian modules. Sharp [8] proved that an injective module
over a Noetherian ring has a secondary representation.

% ============ 习题 ============
\section*{Exercises to Appendix to \S\ref{sec:6}}

\begin{exercise}\label{ex:6.8-ex}
An $A$-module $M$ is \emph{coprimary}\index{coprimary module} if $\Ass(M)$ has just one element. Show that a
finite module $M \neq 0$ over a Noetherian ring $A$ is coprimary if and only if
the following condition is satisfied: for every $a \in A$, the endomorphism
$a\colon M \to M$ is either injective or nilpotent. In this case $\Ass M = \{P\}$,
where $P = \sqrt{\ann M}$.
\end{exercise}

\begin{exercise}\label{ex:6.9-ex}
Show that if $M$ is an $A$-module of finite length then $M$ is coprimary if and
only if it is secondary. Show also that such a module $M$ is a direct sum of
secondary modules belonging to maximal ideals, and $\Ass(M) = \Att(M)$.
\end{exercise}

